\documentclass[a4paper,12pt]{article}

% ------------------------------------------------------------
% Кодировка и язык
% ------------------------------------------------------------
\usepackage[T2A]{fontenc}
\usepackage[utf8]{inputenc}
\usepackage[russian]{babel}

% ------------------------------------------------------------
% Математика
% ------------------------------------------------------------
\usepackage{amsmath,amssymb,amsthm,mathtools}
\usepackage{icomma}

% ------------------------------------------------------------
% Шрифт и типографика
% ------------------------------------------------------------
\usepackage{helvet}
\renewcommand{\familydefault}{\sfdefault}
\usepackage{microtype}
\usepackage{setspace}
\usepackage{parskip}
\usepackage{enumitem}

% ------------------------------------------------------------
% Геометрия страницы
% ------------------------------------------------------------
\usepackage{geometry}
\geometry{
  left=20mm,
  right=20mm,
  top=20mm,
  bottom=22mm
}

% ------------------------------------------------------------
% Цвет
% ------------------------------------------------------------
\usepackage{xcolor}

\definecolor{accent}{HTML}{008080}
\definecolor{darkaccent}{HTML}{004D4D}
\definecolor{textgray}{HTML}{333333}
\definecolor{softgray}{HTML}{707070}
\definecolor{linegray}{HTML}{D8DEDE}

\color{textgray}

% ------------------------------------------------------------
% Графика
% ------------------------------------------------------------
\usepackage{tikz}
\usetikzlibrary{calc}
\usetikzlibrary{arrows.meta,positioning,shapes.geometric}

\usepackage{tkz-euclide}

\usepackage{pgfplots}
\pgfplotsset{compat=1.18}

% ------------------------------------------------------------
% Таблицы
% ------------------------------------------------------------
\usepackage{tabularx,booktabs,longtable}

\newcolumntype{Y}{>{\centering\arraybackslash}X}

% ------------------------------------------------------------
% Колонтитулы и ссылки
% ------------------------------------------------------------
\usepackage{fancyhdr}
\usepackage{hyperref}

\hypersetup{
  colorlinks=true,
  linkcolor=darkaccent,
  urlcolor=darkaccent,
  citecolor=darkaccent
}

\pagestyle{fancy}
\fancyhf{}
\fancyhead[L]{\small ЕГЭ. Профильная математика}
\fancyhead[R]{\small Алгебраическая база и степени}
\fancyfoot[C]{\small\thepage}
\renewcommand{\headrulewidth}{0.3pt}
\renewcommand{\footrulewidth}{0pt}

% ------------------------------------------------------------
% Общие настройки
% ------------------------------------------------------------
\onehalfspacing
\setlength{\parindent}{0pt}
\setlength{\parskip}{0.55em}
\setlist[itemize]{leftmargin=6mm,itemsep=2pt,topsep=3pt}
\setlist[enumerate]{leftmargin=8mm,itemsep=5pt,topsep=4pt}

\emergencystretch=2em

% ------------------------------------------------------------
% Служебные команды
% ------------------------------------------------------------
\newcommand{\MainSection}[1]{%
  \vspace{0.7em}
  {\Large\bfseries\color{darkaccent}#1\par}
  \vspace{0.25em}
  {\color{accent}\rule{\linewidth}{0.6pt}}
  \vspace{0.5em}
}

\newcommand{\SubSection}[1]{%
  \vspace{0.5em}
  {\large\bfseries\color{accent}#1\par}
  \vspace{0.2em}
}

\newcommand{\important}[1]{\textbf{\color{darkaccent}#1}}

\newcommand{\R}{\mathbb{R}}

% ------------------------------------------------------------
% TikZ: единые стили блок-схем
% ------------------------------------------------------------
\tikzset{
  flowstart/.style={
    ellipse,
    draw=darkaccent,
    line width=0.8pt,
    minimum width=42mm,
    minimum height=11mm,
    align=center,
    inner sep=3mm
  },
  flowaction/.style={
    rounded corners=2mm,
    rectangle,
    draw=accent,
    line width=0.8pt,
    text width=98mm,
    minimum height=12mm,
    align=center,
    inner sep=3mm
  },
  flowdecision/.style={
    diamond,
    aspect=2.6,
    draw=darkaccent,
    line width=0.8pt,
    text width=42mm,
    align=center,
    inner sep=1.5mm
  },
  flowside/.style={
    rounded corners=2mm,
    rectangle,
    draw=accent,
    line width=0.8pt,
    text width=46mm,
    minimum height=18mm,
    align=center,
    inner sep=2.5mm
  },
  flowarrow/.style={
    -{Latex[length=2.8mm,width=2mm]},
    line width=0.8pt,
    draw=darkaccent
  }
}

\begin{document}

% ============================================================
% ТИТУЛЬНЫЙ ЛИСТ
% ============================================================

\begin{titlepage}
\thispagestyle{empty}

\begin{center}

\vspace*{2.2cm}

{\small\color{softgray}
ЕГЭ ПО ПРОФИЛЬНОЙ МАТЕМАТИКЕ
}

\vspace{1.3cm}

{\Huge\bfseries\color{darkaccent}
Чек-лист
}

\vspace{0.8cm}

{\LARGE\bfseries
Алгебраическая база\\[0.25em]
и степени
}

\vspace{0.55cm}

{\Large\color{accent}
Приведение выражений к степенному виду
}

\vspace{1.5cm}

{\large
Формулы сокращённого умножения
}

\vspace{0.3cm}

{\large
Квадратные выражения, уравнения и неравенства
}

\vspace{0.3cm}

{\large
Свойства степеней
}

\vspace{0.3cm}

{\large
Составные числа, корни и дроби
}

\vfill

{\large
Матвей Горбачёв
}

\vspace{0.8cm}

{\normalsize
\today
}

\vspace{1.0cm}

{\small
Лёвин Артём Александрович\\
эксклюзивно для Матвея Горбачёва
}

\vspace{1.4cm}

\begin{tikzpicture}
\draw[
  accent,
  line width=1.1pt
]
(-6,0)
.. controls (-3.8,0.8) and (-2.3,-0.65) ..
(0,0)
.. controls (2.2,0.65) and (4.0,-0.75) ..
(6,0);
\end{tikzpicture}

\end{center}

\end{titlepage}

% ============================================================
% 1. АЛГЕБРАИЧЕСКАЯ РАЗМИНКА
% ============================================================

\MainSection{1. Алгебраическая разминка}

Цель этого блока --- быстро восстановить инструменты, которые постоянно
используются далее в степенях, логарифмах, производной и других темах ЕГЭ.

\SubSection{1.1. Вынесение общего множителя}

Если несколько слагаемых имеют общий множитель, его можно вынести за скобки:

\[
3x+6y=3(x+2y).
\]

Полезно всегда проверять результат обратным действием:

\[
3(x+2y)=3x+6y.
\]

\important{Принцип:} преобразование должно быть обратимым и проверяемым.

% ------------------------------------------------------------

\SubSection{1.2. Квадрат суммы и квадрат разности}

\[
(a+b)^2=a^2+2ab+b^2,
\]

\[
(a-b)^2=a^2-2ab+b^2.
\]

Пример:

\[
(2x+3)^2
=
(2x)^2+2\cdot 2x\cdot 3+3^2
=
4x^2+12x+9.
\]

Обратите внимание: удваивается \important{произведение} первого и второго
элементов.

% ------------------------------------------------------------

\SubSection{1.3. Разность квадратов}

\[
a^2-b^2=(a-b)(a+b).
\]

Например,

\[
4k^2-9t^2
=
(2k)^2-(3t)^2
=
(2k-3t)(2k+3t).
\]

Формула работает и справа налево:

\[
(2k-3t)(2k+3t)=4k^2-9t^2.
\]

% ------------------------------------------------------------

\SubSection{1.4. Куб суммы и куб разности}

\[
(a+b)^3
=
a^3+3a^2b+3ab^2+b^3,
\]

\[
(a-b)^3
=
a^3-3a^2b+3ab^2-b^3.
\]

Например,

\[
(x+1)^3=x^3+3x^2+3x+1.
\]

При кубе разности знаки идут по схеме

\[
+\quad -\quad +\quad -.
\]

% ============================================================
% 2. КВАДРАТНЫЙ ТРЁХЧЛЕН
% ============================================================

\MainSection{2. Квадратный трёхчлен}

Пусть

\[
ax^2+bx+c,
\qquad a\ne 0.
\]

Если квадратное уравнение

\[
ax^2+bx+c=0
\]

имеет действительные корни \(x_1\) и \(x_2\), то

\[
ax^2+bx+c
=
a(x-x_1)(x-x_2).
\]

\SubSection{Пример}

Разложим на множители

\[
x^2-x-6.
\]

Решаем уравнение

\[
x^2-x-6=0.
\]

Дискриминант:

\[
D=(-1)^2-4\cdot1\cdot(-6)=1+24=25.
\]

Корни:

\[
x_1=\frac{1-5}{2}=-2,
\qquad
x_2=\frac{1+5}{2}=3.
\]

Следовательно,

\[
x^2-x-6
=
(x+2)(x-3).
\]

\SubSection{Теорема Виета}

Для уравнения

\[
ax^2+bx+c=0
\]

выполняется

\[
x_1+x_2=-\frac ba,
\qquad
x_1x_2=\frac ca.
\]

В приведённом уравнении

\[
x^2+bx+c=0
\]

получаем

\[
x_1+x_2=-b,
\qquad
x_1x_2=c.
\]

% ============================================================
% 3. БАЗОВЫЕ УРАВНЕНИЯ
% ============================================================

\MainSection{3. Базовые уравнения}

\SubSection{3.1. Вынесение переменной}

\[
x^2-6x=0.
\]

Выносим \(x\):

\[
x(x-6)=0.
\]

Произведение равно нулю, если хотя бы один множитель равен нулю:

\[
x=0
\qquad\text{или}\qquad
x=6.
\]

Ответ:

\[
\boxed{x\in\{0;6\}}.
\]

% ------------------------------------------------------------

\SubSection{3.2. Уравнение \(x^2=a\)}

Если

\[
x^2=6,
\]

то

\[
x=\pm\sqrt6.
\]

Поэтому

\[
x^2-6=0
\quad\Longrightarrow\quad
x=\pm\sqrt6.
\]

Если же

\[
x^2+6=0,
\]

то

\[
x^2=-6.
\]

В действительных числах такое равенство невозможно:

\[
\boxed{\emptyset}.
\]

\important{Ключевая проверка:}
для любого \(x\in\R\)

\[
x^2\ge 0.
\]

% ============================================================
% 4. НЕРАВЕНСТВА
% ============================================================

\MainSection{4. Неравенства}

\SubSection{4.1. Линейное неравенство}

\[
x+5\ge 0.
\]

Переносим \(5\):

\[
x\ge -5.
\]

Ответ:

\[
\boxed{[-5;+\infty)}.
\]

При нестрогом знаке \(\ge\) граничная точка входит в ответ.

% ------------------------------------------------------------

\SubSection{4.2. Квадратное неравенство}

Решим

\[
x^2-4\ge 0.
\]

Разложим левую часть:

\[
x^2-4=(x-2)(x+2).
\]

Нули:

\[
x=-2,
\qquad
x=2.
\]

На промежутках

\[
(-\infty;-2),\qquad
(-2;2),\qquad
(2;+\infty)
\]

знаки выражения имеют вид

\[
+\qquad -\qquad +.
\]

Следовательно,

\[
\boxed{
x\in(-\infty;-2]\cup[2;+\infty)
}.
\]

\begin{center}
\begin{tikzpicture}[x=1.15cm,y=1cm]

\draw[-{Latex[length=2.5mm]},line width=0.7pt]
(-5,0)--(5.3,0);

\draw (-2,0.10)--(-2,-0.10)
node[below=3pt] {$-2$};

\draw (2,0.10)--(2,-0.10)
node[below=3pt] {$2$};

\fill[accent] (-2,0) circle (2.2pt);
\fill[accent] (2,0) circle (2.2pt);

\draw[
  accent,
  line width=1.4pt,
  -{Latex[length=2.5mm]}
]
(-2,0.28)--(-4.65,0.28);

\draw[
  accent,
  line width=1.4pt,
  -{Latex[length=2.5mm]}
]
(2,0.28)--(4.65,0.28);

\end{tikzpicture}
\end{center}

\SubSection{Когда знаки чередуются?}

В методе интервалов знак выражения меняется при переходе через корень
\important{нечётной кратности}.

Например,

\[
(x-1)(x+3)
\]

имеет простые корни, поэтому при прохождении каждого корня знак меняется.

Если множитель имеет чётную степень,

\[
(x-1)^2,
\]

то при переходе через \(x=1\) знак \important{сохраняется}.

% ============================================================
% 5. СТЕПЕНИ
% ============================================================

\MainSection{5. Основные свойства степеней}

Это один из фундаментальных блоков подготовки к профильному ЕГЭ.
Свойства степеней далее используются в логарифмах, производной,
показательных уравнениях и преобразованиях выражений.

\SubSection{Правило 1. Произведение степеней с одинаковым основанием}

\[
a^m\cdot a^n=a^{m+n}.
\]

Например,

\[
5^3\cdot5^4=5^7.
\]

\important{Условие:} основания должны быть одинаковыми.

Формула

\[
2^3\cdot5^4=10^7
\]

неверна.

% ------------------------------------------------------------

\SubSection{Правило 2. Деление степеней с одинаковым основанием}

При \(a\ne0\):

\[
\frac{a^m}{a^n}=a^{m-n}.
\]

Например,

\[
\frac{5^4}{5^3}=5.
\]

Отсюда также

\[
a^0=1,
\qquad a\ne0.
\]

% ------------------------------------------------------------

\SubSection{Правило 3. Степень произведения}

\[
(ab)^n=a^nb^n.
\]

Формула работает и справа налево:

\[
a^nb^n=(ab)^n.
\]

Например,

\[
42^3
=
(6\cdot7)^3
=
6^3\cdot7^3.
\]

\important{Важно:}

\[
(a+b)^n
\]

нельзя преобразовывать в

\[
a^n+b^n.
\]

% ------------------------------------------------------------

\SubSection{Правило 4. Степень частного}

При \(b\ne0\):

\[
\left(\frac ab\right)^n
=
\frac{a^n}{b^n}.
\]

И обратно:

\[
\frac{a^n}{b^n}
=
\left(\frac ab\right)^n.
\]

% ------------------------------------------------------------

\SubSection{Правило 5. Степень степени}

\[
(a^m)^n=a^{mn}.
\]

Например,

\[
(5^3)^2=5^6.
\]

Сравните:

\[
(5^3)^2=5^6,
\]

но

\[
5^3\cdot5^2=5^5.
\]

В первом случае показатели \important{перемножаются},
во втором --- \important{складываются}.

Кроме того,

\[
(a^m)^n=(a^n)^m.
\]

% ------------------------------------------------------------

\SubSection{Правило 6. Отрицательная степень}

При \(a\ne0\):

\[
\frac1{a^n}=a^{-n}.
\]

Например,

\[
\frac1{5^3}=5^{-3}.
\]

Если явный показатель отсутствует,

\[
\frac15=5^{-1}.
\]

Полезное следствие:

\[
\left(\frac ab\right)^n
=
\left(\frac ba\right)^{-n},
\qquad
a\ne0,\quad b\ne0.
\]

Например,

\[
\frac23
=
\left(\frac32\right)^{-1}.
\]

% ------------------------------------------------------------

\SubSection{Правило 7. Корень как степень}

Для допустимых значений выражения:

\[
\sqrt[k]{a^m}
=
a^{\frac{m}{k}}.
\]

Например,

\[
\sqrt[5]{2^3}=2^{3/5},
\]

а

\[
\sqrt2=2^{1/2}.
\]

Мнемоническая идея:

\[
\boxed{
\text{внутренняя степень}
\div
\text{внешняя степень}
}
\]

то есть показатель внутри радикала становится числителем,
а показатель корня --- знаменателем.

% ============================================================
% 6. ОДЗ И ГРАНИЦЫ ПРИМЕНИМОСТИ
% ============================================================

\MainSection{6. ОДЗ и границы применимости}

При преобразованиях степеней важно учитывать область определения.

\begin{itemize}

\item
В дроби знаменатель не может быть равен нулю.

\[
\frac1{a^n}
\quad\Longrightarrow\quad
a\ne0.
\]

\item
Для корня чётной степени требуется

\[
\sqrt[2k]{A}
\quad\Longrightarrow\quad
A\ge0.
\]

\item
Если корень чётной степени находится в знаменателе, требуется

\[
A>0.
\]

\item
Корень нечётной степени определён для любого действительного
подкоренного выражения:

\[
\sqrt[3]{-8}=-2.
\]

\item
При работе с произвольными действительными показателями
основание обычно рассматривается положительным:

\[
a^\alpha,
\qquad a>0.
\]

Это особенно важно в показательных и логарифмических задачах ЕГЭ.

\end{itemize}

% ============================================================
% 7. СОСТАВНЫЕ И ПРОСТЫЕ ЧИСЛА
% ============================================================

\MainSection{7. Простые и составные числа}

\SubSection{Простое число}

Натуральное число \(p>1\) называется простым, если имеет ровно два
положительных делителя:

\[
1
\quad\text{и}\quad
p.
\]

Примеры:

\[
2,\quad3,\quad5,\quad7,\quad11,\quad13.
\]

\SubSection{Составное число}

Натуральное число называется составным, если оно имеет более двух
положительных делителей.

Например,

\[
8=2^3,
\]

\[
54=2\cdot27
=2\cdot3^3.
\]

\important{Число \(1\)} не является ни простым, ни составным.

% ============================================================
% 8. ГЛАВНАЯ СТРАТЕГИЯ
% ============================================================

\MainSection{8. Главная стратегия задач на степени}

Основной принцип занятия:

\[
\boxed{
\text{Приводим всё к степенному виду}
}
\]

В первую очередь ищем три типа объектов:

\[
\boxed{
\text{составные числа}
\quad
\text{корни}
\quad
\text{дроби}
}
\]

После этого стараемся привести выражение к одинаковым основаниям
и применить свойства степеней.

\SubSection{Путь 1. Составные числа}

Разлагаем их на простые множители.

Например,

\[
81=3^4,
\]

\[
25=5^2,
\]

\[
54=2\cdot3^3.
\]

\SubSection{Путь 2. Корни}

Переводим радикалы в дробные показатели:

\[
\sqrt[3]{81}
=
81^{1/3}
=
(3^4)^{1/3}
=
3^{4/3}.
\]

\SubSection{Путь 3. Дроби}

Используем отрицательные степени:

\[
\frac1{81}
=
\frac1{3^4}
=
3^{-4}.
\]

Если встречается несколько типов объектов, применяем преобразования
последовательно.

% ============================================================
% БЛОК-СХЕМА
% ============================================================

\clearpage

\begin{center}

{\LARGE\bfseries\color{darkaccent}
Алгоритм: как решать выражения со степенями
}

\vspace{1.1cm}

\begin{tikzpicture}[
  node distance=8mm and 13mm
]

\node[flowstart] (start)
{Получено выражение};

\node[flowaction,below=of start] (scan)
{Просмотрите выражение целиком:
найдите составные числа, корни, дроби,
разные основания и сложные показатели};

\node[flowdecision,below=10mm of scan] (ready)
{Все основные элементы уже записаны в степенном виде?};

\node[flowside,left=16mm of ready] (convert)
{Преобразуйте:\\
составные числа --- разложить;\\
корни --- дробные степени;\\
дроби --- отрицательные степени};

\node[flowaction,below=11mm of ready] (bases)
{По возможности приведите степени к одинаковым основаниям};

\node[flowaction,below=of bases] (rules)
{Выберите нужное свойство:
сложение или вычитание показателей,
умножение показателей,
вынесение общей степени};

\node[flowaction,below=of rules] (simplify)
{Аккуратно выполните действия с показателями.
Вычитаемые выражения предварительно заключайте в скобки};

\node[flowaction,below=of simplify] (check)
{Проверьте ОДЗ, знаменатели, корни,
знаки и допустимость каждого преобразования};

\node[flowstart,below=of check] (finish)
{Запишите окончательный ответ};

\draw[flowarrow] (start)--(scan);
\draw[flowarrow] (scan)--(ready);

\draw[flowarrow]
(ready.west)
-- node[above,pos=0.42]{\small нет}
(convert.east);

\draw[flowarrow]
(convert.north)
|- ($(scan.west)+(0,-0.15)$);

\draw[flowarrow]
(ready)
-- node[right,pos=0.38]{\small да}
(bases);

\draw[flowarrow] (bases)--(rules);
\draw[flowarrow] (rules)--(simplify);
\draw[flowarrow] (simplify)--(check);
\draw[flowarrow] (check)--(finish);

\end{tikzpicture}

\end{center}

\clearpage

% ============================================================
% 9. ТИПОВЫЕ ПРИМЕРЫ
% ============================================================

\MainSection{9. Типовые задачи}

\SubSection{Пример 1. Десятичные показатели}

Вычислим:

\[
5^{0{,}36}\cdot25^{0{,}32}.
\]

Число \(25\) --- составное:

\[
25=5^2.
\]

Подставляем:

\[
5^{0{,}36}\cdot(5^2)^{0{,}32}.
\]

Степень степени:

\[
(5^2)^{0{,}32}=5^{0{,}64}.
\]

Следовательно,

\[
5^{0{,}36}\cdot5^{0{,}64}
=
5^{0{,}36+0{,}64}
=
5^1
=
\boxed5.
\]

% ------------------------------------------------------------

\SubSection{Пример 2. Несколько оснований}

Вычислим:

\[
\frac{
2^{3{,}5}\cdot3^{5{,}5}
}{
6^{4{,}5}
}.
\]

Раскладываем

\[
6=2\cdot3.
\]

Тогда

\[
6^{4{,}5}
=
(2\cdot3)^{4{,}5}
=
2^{4{,}5}3^{4{,}5}.
\]

Получаем

\[
\frac{
2^{3{,}5}3^{5{,}5}
}{
2^{4{,}5}3^{4{,}5}
}
=
2^{3{,}5-4{,}5}
3^{5{,}5-4{,}5}.
\]

Следовательно,

\[
2^{-1}3^1
=
\frac32.
\]

Ответ:

\[
\boxed{\frac32}.
\]

% ------------------------------------------------------------

\SubSection{Пример 3. Корни разных степеней}

Вычислим:

\[
\left(
\frac{
2^{1/3}\sqrt[4]{2}
}{
\sqrt[12]{2}
}
\right)^2.
\]

Переводим корни в степени:

\[
\sqrt[4]{2}=2^{1/4},
\qquad
\sqrt[12]{2}=2^{1/12}.
\]

Тогда

\[
\left(
2^{1/3+1/4-1/12}
\right)^2.
\]

Приводим показатели к знаменателю \(12\):

\[
\frac13+\frac14-\frac1{12}
=
\frac4{12}+\frac3{12}-\frac1{12}
=
\frac6{12}
=
\frac12.
\]

Поэтому

\[
\left(2^{1/2}\right)^2
=
2.
\]

Ответ:

\[
\boxed2.
\]

% ------------------------------------------------------------

\SubSection{Пример 4. Показатели с радикалами}

Рассмотрим

\[
\frac{
5^{3\sqrt7-1}\cdot
5^{1-\sqrt7}
}{
5^{2\sqrt7-1}
}.
\]

Основания одинаковые, поэтому работаем только с показателями:

\[
5^{
(3\sqrt7-1)
+
(1-\sqrt7)
-
(2\sqrt7-1)
}.
\]

Раскрываем скобки:

\[
3\sqrt7-1+1-\sqrt7-2\sqrt7+1=1.
\]

Следовательно,

\[
\boxed5.
\]

\important{Критически важный приём:}
если из показателя вычитается сложное выражение,
его сначала заключают в скобки.

% ============================================================
% 10. ПРАВИЛО СКОБОК
% ============================================================

\MainSection{10. Скобки как защита от ошибок}

Скобки особенно важны в трёх ситуациях.

\begin{enumerate}

\item
\important{Подстановка сложного выражения.}

Если

\[
a=2x+3,
\]

то в выражение \(a^2\) подставляем

\[
(2x+3)^2.
\]

\item
\important{Вычитание сложного выражения.}

Пишем

\[
A-(B+C),
\]

а затем раскрываем:

\[
A-B-C.
\]

\item
\important{Возведение произведения или суммы в степень.}

\[
(2\cdot3)^n,
\qquad
(x+2)^2.
\]

\end{enumerate}

Полезная привычка:

\[
\boxed{
\text{сначала поставить скобки,
затем выполнять преобразование}
}
\]

Это занимает несколько секунд и резко снижает число знаковых ошибок.

% ============================================================
% 11. ТИПИЧНЫЕ ОШИБКИ
% ============================================================

\MainSection{11. Типичные ошибки}

\begin{enumerate}

\item
Складывать показатели при разных основаниях:

\[
2^3\cdot3^4
\ne
6^7.
\]

\item
Распространять степень на сумму:

\[
(a+b)^2
\ne
a^2+b^2.
\]

Верно:

\[
(a+b)^2=a^2+2ab+b^2.
\]

\item
Путать степень степени с произведением степеней:

\[
(a^m)^n=a^{mn},
\]

но

\[
a^m\cdot a^n=a^{m+n}.
\]

\item
Забывать минус при обратной дроби:

\[
\frac ab
=
\left(\frac ba\right)^{-1}.
\]

\item
Ошибаться при вычитании сложного показателя:

\[
a^{m-n}
\]

требует аккуратной работы со всем выражением \(n\).

\item
Игнорировать ОДЗ дробей и корней.

\item
Оставлять составное число в неудобном виде, когда его можно
представить через уже имеющееся основание.

Например,

\[
25=5^2,
\qquad
27=3^3,
\qquad
81=3^4.
\]

\item
При методе интервалов автоматически чередовать знаки,
не проверив кратность корней.

\end{enumerate}

% ============================================================
% 12. КРАТКИЙ ЧЕК-ЛИСТ
% ============================================================

\MainSection{12. Чек-лист перед записью ответа}

Перед завершением задачи проверьте:

\begin{itemize}

\item
Все ли составные числа удалось разложить удобным способом?

\item
Можно ли заменить корни дробными степенями?

\item
Можно ли заменить дроби отрицательными степенями?

\item
Можно ли получить одинаковые основания?

\item
Какое свойство степени применяется именно здесь?

\item
Нужны ли скобки при подстановке или вычитании?

\item
Корректно ли выполнены действия с дробными показателями?

\item
Проверены ли ограничения:
знаменатель, корень чётной степени, нулевое основание?

\item
Можно ли быстро проверить полученный результат обратным преобразованием?

\end{itemize}

% ============================================================
% ТРЕНИРОВОЧНЫЙ БЛОК
% ============================================================

\clearpage

\begin{center}
{\LARGE\bfseries\color{darkaccent}
Тренировочный блок
}

\vspace{0.25cm}

{\large
Алгебраическая база и степени
}

\vspace{0.2cm}

{\small\color{softgray}
Путь А
}
\end{center}

\vspace{0.5cm}

Задания расположены по нарастающей сложности.
Решения записывайте полностью: преобразование оснований,
работу с показателями и необходимые ограничения.

\SubSection{Средний уровень}

\begin{enumerate}[label=\textbf{\arabic*.}]

\item
Упростите выражение

\[
(2x+3)^2-(2x-3)^2.
\]

\item
Разложите на множители:

\[
x^2-5x-14.
\]

\item
Решите неравенство:

\[
x^2-16\ge0.
\]

\end{enumerate}

\SubSection{Уровень выше среднего}

\begin{enumerate}[label=\textbf{\arabic*.},start=4]

\item
Вычислите:

\[
7^{0{,}28}\cdot49^{0{,}36}.
\]

\item
Вычислите:

\[
\frac{
2^{5/2}\cdot3^{7/2}
}{
6^{5/2}
}.
\]

\item
Вычислите:

\[
\left(
\frac{
2^{1/3}\sqrt[4]{2}
}{
\sqrt[12]{2}
}
\right)^2.
\]

\item
Вычислите:

\[
\left(\frac35\right)^{-2}
\cdot
\frac9{25}.
\]

\item
Вычислите:

\[
\frac{
5^{2\sqrt3-1}\cdot5^{4-\sqrt3}
}{
5^{\sqrt3+2}
}.
\]

\item
Вычислите:

\[
\frac{\sqrt[3]{54}}{\sqrt[3]{2}}.
\]

\end{enumerate}

\SubSection{Сложный уровень}

\begin{enumerate}[label=\textbf{\arabic*.},start=10]

\item
Вычислите:

\[
\frac{
81^{3/4}\cdot27^{-1/3}
}{
3^{-2}
}
\cdot
\left(
\frac{
\sqrt[6]{9}
}{
\sqrt[4]{3}
}
\right)^{12}.
\]

\end{enumerate}

% ============================================================
% ОТВЕТЫ
% ============================================================

\clearpage

\begin{center}

{\LARGE\bfseries\color{darkaccent}
Ответы
}

\vspace{0.4cm}

{\small\color{softgray}
Сначала выполните задания самостоятельно.
}

\end{center}

\vspace{0.8cm}

\renewcommand{\arraystretch}{1.45}

\begin{center}
\begin{tabularx}{\linewidth}{@{}cY@{}}
\toprule
\textbf{№} & \textbf{Ответ} \\
\midrule

1 &
\[
24x
\]
\\

2 &
\[
(x-7)(x+2)
\]
\\

3 &
\[
(-\infty;-4]\cup[4;+\infty)
\]
\\

4 &
\[
7
\]
\\

5 &
\[
3
\]
\\

6 &
\[
2
\]
\\

7 &
\[
1
\]
\\

8 &
\[
5
\]
\\

9 &
\[
3
\]
\\

10 &
\[
243
\]
\\

\bottomrule
\end{tabularx}
\end{center}

\vfill

\begin{center}
{\color{accent}\rule{0.72\linewidth}{0.5pt}}

\vspace{0.45cm}

{\small\color{softgray}
Главная идея занятия:
}

\vspace{0.2cm}

{\large\bfseries\color{darkaccent}
привести выражение к удобному степенному виду,
а затем применить свойства степеней.
}

\end{center}

\end{document}